% !TEX root = GRE-Textbook.tex

\documentclass[openany]{tufte-book}

\usepackage{amsfonts, amsmath, amsthm, amssymb}
\usepackage{adjustbox, caption, enumerate, fancyhdr, float, graphicx, hyperref, mathtools, minitoc, multicol, tasks, thmtools, tikz, titlesec, venndiagram, xcolor}
\usepackage[table]{xcolor}
%\usepackage{sectsty}
\usepackage[framemethod=TikZ]{mdframed}

\definecolor{fsblue}{RGB}{132,175,230}
\declaretheoremstyle[  % Blue theorem boxes
	headfont=\bfseries\color{fsblue!60!black},
	bodyfont=\normalfont,
	mdframed={
     	linewidth=2pt,
     	rightline=false, topline=false, bottomline=false,
     	linecolor=fsblue, backgroundcolor=fsblue!10,
    }]{thmblue}
    
\declaretheoremstyle[  % Blue lemma/corollary boxes
	headfont=\bfseries\color{fsblue!60!black},
	bodyfont=\normalfont,
	mdframed={
        	linewidth=2pt,
        	rightline=false, topline=false, bottomline=false,
        	linecolor=fsblue,
    }]{lemblue}
    
\declaretheoremstyle[  % Blue proofs
	headfont=\itshape\color{fsblue!60!black},
	notefont = \itshape ,
	name = Proof,
	numbered = no,
	bodyfont=\normalfont,
	qed = \qedsymbol,
	mdframed={
        	linewidth=2pt,
        	rightline=false, topline=false, bottomline=false,
        	linecolor=fsblue,
    }]{pfblue}
    
\declaretheoremstyle[  % Box for definitions
	headfont=\bfseries,
	bodyfont=\normalfont,
	numbered=no,
	mdframed={
        	linewidth=1.5pt,
        	rightline=false, topline=false, bottomline=false,
    }]{defbox}
    
\declaretheoremstyle[  % Box for problems
	headfont=\bfseries,
	bodyfont=\normalfont,
	mdframed={
        	linewidth=1.5pt,
        	rightline=false, topline=false, bottomline=false,
    }]{probbox}

\newenvironment{multienum}[1]{\begin{multicols}{#1} \begin{enumerate}[$i$.)]}{\end{enumerate}\end{multicols}}

% asmthm Environments
\theoremstyle{plain}
\declaretheorem[style=thmblue, parent=chapter]{theorem}  % Theorems
\declaretheorem[style=thmblue, sibling=theorem]{proposition}  % Propositions
\declaretheorem[style=lemblue, parent=theorem]{lemma}  % Lemmas
\declaretheorem[style=lemblue, parent=theorem]{corollary}  % Corollaries

\theoremstyle{definition}
\declaretheorem[style=defbox]{definition}  % Definitions
\declaretheorem[numbered=no]{example}  % Examples
\declaretheorem[style=probbox]{problem}  % Problems
\declaretheorem[numbered=no]{remark}  % Remarks
\newenvironment{prob}[1]  % Problem with custom number
  {\renewcommand\theproblem{#1}\problem} {\endproblem}

\theoremstyle{remark}
\newtheorem{step}{Step}  % Steps
\newtheorem{case}{Case}  % Cases
\newtheorem{Rule}{Rule}  % Rules
\declaretheorem[style=pfblue]{Prooff}  % Proofs in color
	\newenvironment{Proof}[1][ ]{\vspace{-1.75em}\begin{Prooff}[#1]}{\end{Prooff}}  % Numbered proof, removes space between thm and pf box
\declaretheorem[numbered=no]{solution}  % Solutions

% Shortcuts for number systems
\newcommand{\N}{\ensuremath{\mathbb{N}}}
\newcommand{\Z}{\ensuremath{\mathbb{Z}}}
\newcommand{\Q}{\ensuremath{\mathbb{Q}}}
\newcommand{\R}{\ensuremath{\mathbb{R}}}
\newcommand{\C}{\ensuremath{\mathbb{C}}}

% Operators
\DeclareMathOperator{\Ker}{Ker}  % Kernel
\DeclareMathOperator{\dom}{dom}  % Domain
\DeclareMathOperator{\Arg}{Arg}  % Principal argument of complex num
\DeclareMathOperator{\Log}{Log}  % Complex logarithm
\DeclareMathOperator{\conv}{conv}  % Convex hull
\DeclareMathOperator{\Hom}{Hom}  % Morphisms

\newcommand{\be}{\mathbf{e}}
\newcommand{\mcA}{\mathcal{A}}
\newcommand{\mcE}{\mathcal{E}}
\newcommand{\mcM}{\mathcal{M}}
\newcommand{\mcR}{\mathcal{R}}
\newcommand{\borR}{\mathcal{B}_\R}
\newcommand{\empt}{\varnothing}
\newcommand{\llra}{\longleftrightarrow}


% Avoids widows and orphans
\clubpenalty = 10000
\widowpenalty = 10000
\displaywidowpenalty = 10000


\hypersetup{colorlinks}% uncomment this line if you prefer colored hyperlinks (e.g., for onscreen viewing)

\fancypagestyle{plain}{}

%% Defines a subtitle command
\makeatletter
\newcommand{\plainsubtitle}{}%     plain-text-only subtitle
\newcommand{\subtitle}[1]{%
  \gdef\@subtitle{#1}%
  \renewcommand{\plainsubtitle}{#1}% use provided plain-text title
  \ifthenelse{\isundefined{\hypersetup}}%
    {}% hyperref is not loaded; do nothing
    {\hypersetup{pdftitle={\plaintitle: \plainsubtitle{}}}}% set the PDF metadata title
}

%% Makes title page resemble Tufte's VE instead of default
\renewcommand{\maketitlepage}{%
  {%
  \begin{fullwidth}%
  \fontsize{20}{20}\selectfont\par\noindent\textcolor{darkgray}{\thanklessauthor}%
  \vspace{11.5pc}%
  \fontsize{35}{40}\selectfont\par\noindent\textcolor{darkgray}{\thanklesstitle}%
  \fontsize{20}{40}\selectfont\par\noindent\textcolor{darkgray}{\plainsubtitle}%
  \vfill
  \fontsize{14}{16}\selectfont\par\noindent\thanklesspublisher%
  \end{fullwidth}%
  }%
  \thispagestyle{empty}%
}

%% Makes TOC resemble Tufte's VDQI instead of default
\titlecontents{part}%
    [0pt]% distance from left margin
    {\addvspace{0.25\baselineskip}}% above (global formatting of entry)
    {\allcaps{Part~\thecontentslabel}\allcaps}% before w/ label (label = ``Part I'')
    {\allcaps{Part~\thecontentslabel}\allcaps}% before w/o label
    {}% filler and page (leaders and page num)
    [\vspace*{0.5\baselineskip}]% after
\titlecontents{chapter}%
    [4em]% distance from left margin
    {}% above (global formatting of entry)
    {\contentslabel{2em}\large}% before w/ label (label = ``Chapter 1'')
    {\hspace{0em}\large}% before w/o label
    {\dotfill\thecontentspage}% filler and page (leaders and page num)
    [\vspace*{0.25\baselineskip}]% after
\titlecontents{section}%
    [8em]% distance from left margin
    {}% above (global formatting of entry)
    {\contentslabel{3em}\textit}% before w/ label (label = ``Section 1.1'')
    {\hspace{0em}\textit}% before w/o label
    {\qquad\thecontentspage}% filler and page (leaders and page num)
    [\vspace*{0.25\baselineskip}]% after

%%
% If they're installed, use Bergamo and Chantilly from www.fontsite.com.
% They're clones of Bembo and Gill Sans, respectively.
\IfFileExists{bergamo.sty}{\usepackage[osf]{bergamo}}{}% Bembo
\IfFileExists{chantill.sty}{\usepackage{chantill}}{}% Gill Sans

%\usepackage{microtype}

% Just some sample text
\usepackage{lipsum}

% For nicely typeset tabular material
\usepackage{booktabs}

% For graphics / images
\usepackage{graphicx}
\setkeys{Gin}{width=\linewidth,totalheight=\textheight,keepaspectratio}
\graphicspath{{graphics/}}

% The fancyvrb package lets us customize the formatting of verbatim
% environments.  We use a slightly smaller font.
\usepackage{fancyvrb}
\fvset{fontsize=\normalsize}

% Prints argument within hanging parentheses (i.e., parentheses that take
% up no horizontal space).  Useful in tabular environments.
\newcommand{\hangp}[1]{\makebox[0pt][r]{(}#1\makebox[0pt][l]{)}}

% Prints an asterisk that takes up no horizontal space.
% Useful in tabular environments.
\newcommand{\hangstar}{\makebox[0pt][l]{*}}

% Prints a trailing space in a smart way.
\usepackage{xspace}

% Inserts a blank page
\newcommand{\blankpage}{\newpage\hbox{}\thispagestyle{empty}\newpage}

\usepackage{units}

% Typesets the font size, leading, and measure in the form of 10/12x26 pc.
\newcommand{\measure}[3]{#1/#2$\times$\unit[#3]{pc}}

% Macros for typesetting the documentation
\newcommand{\hlred}[1]{\textcolor{Maroon}{#1}}% prints in red
\newcommand{\hangleft}[1]{\makebox[0pt][r]{#1}}
\newcommand{\hairsp}{\hspace{1pt}}% hair space
\newcommand{\hquad}{\hskip0.5em\relax}% half quad space
\newcommand{\TODO}{\textcolor{red}{\bf TODO!}\xspace}
\newcommand{\ie}{\textit{i.\hairsp{}e.}\xspace}
\newcommand{\eg}{\textit{e.\hairsp{}g.}\xspace}
\newcommand{\na}{\quad--}% used in tables for N/A cells
\providecommand{\XeLaTeX}{X\lower.5ex\hbox{\kern-0.15em\reflectbox{E}}\kern-0.1em\LaTeX}
\newcommand{\tXeLaTeX}{\XeLaTeX\index{XeLaTeX@\protect\XeLaTeX}}
% \index{\texttt{\textbackslash xyz}@\hangleft{\texttt{\textbackslash}}\texttt{xyz}}
\newcommand{\tuftebs}{\symbol{'134}}% a backslash in tt type in OT1/T1
\newcommand{\doccmdnoindex}[2][]{\texttt{\tuftebs#2}}% command name -- adds backslash automatically (and doesn't add cmd to the index)
\newcommand{\doccmddef}[2][]{%
  \hlred{\texttt{\tuftebs#2}}\label{cmd:#2}%
  \ifthenelse{\isempty{#1}}%
    {% add the command to the index
      \index{#2 command@\protect\hangleft{\texttt{\tuftebs}}\texttt{#2}}% command name
    }%
    {% add the command and package to the index
      \index{#2 command@\protect\hangleft{\texttt{\tuftebs}}\texttt{#2} (\texttt{#1} package)}% command name
      \index{#1 package@\texttt{#1} package}\index{packages!#1@\texttt{#1}}% package name
    }%
}% command name -- adds backslash automatically
\newcommand{\doccmd}[2][]{%
  \texttt{\tuftebs#2}%
  \ifthenelse{\isempty{#1}}%
    {% add the command to the index
      \index{#2 command@\protect\hangleft{\texttt{\tuftebs}}\texttt{#2}}% command name
    }%
    {% add the command and package to the index
      \index{#2 command@\protect\hangleft{\texttt{\tuftebs}}\texttt{#2} (\texttt{#1} package)}% command name
      \index{#1 package@\texttt{#1} package}\index{packages!#1@\texttt{#1}}% package name
    }%
}% command name -- adds backslash automatically
\newcommand{\docopt}[1]{\ensuremath{\langle}\textrm{\textit{#1}}\ensuremath{\rangle}}% optional command argument
\newcommand{\docarg}[1]{\textrm{\textit{#1}}}% (required) command argument
\newenvironment{docspec}{\begin{quotation}\ttfamily\parskip0pt\parindent0pt\ignorespaces}{\end{quotation}}% command specification environment
\newcommand{\docenv}[1]{\texttt{#1}\index{#1 environment@\texttt{#1} environment}\index{environments!#1@\texttt{#1}}}% environment name
\newcommand{\docenvdef}[1]{\hlred{\texttt{#1}}\label{env:#1}\index{#1 environment@\texttt{#1} environment}\index{environments!#1@\texttt{#1}}}% environment name
\newcommand{\docpkg}[1]{\texttt{#1}\index{#1 package@\texttt{#1} package}\index{packages!#1@\texttt{#1}}}% package name
\newcommand{\doccls}[1]{\texttt{#1}}% document class name
\newcommand{\docclsopt}[1]{\texttt{#1}\index{#1 class option@\texttt{#1} class option}\index{class options!#1@\texttt{#1}}}% document class option name
\newcommand{\docclsoptdef}[1]{\hlred{\texttt{#1}}\label{clsopt:#1}\index{#1 class option@\texttt{#1} class option}\index{class options!#1@\texttt{#1}}}% document class option name defined
\newcommand{\docmsg}[2]{\bigskip\begin{fullwidth}\noindent\ttfamily#1\end{fullwidth}\medskip\par\noindent#2}
\newcommand{\docfilehook}[2]{\texttt{#1}\index{file hooks!#2}\index{#1@\texttt{#1}}}
\newcommand{\doccounter}[1]{\texttt{#1}\index{#1 counter@\texttt{#1} counter}}

% Generates the index
\usepackage{makeidx}
\makeindex

% Metadata
\title{Fundamentals of Mathematics}
\subtitle{for the GRE Quantitative Reasoning Section}
\author{\itshape Carson Mitchell}
\publisher{}

\setcounter{secnumdepth}{3}